\documentclass[12pt]{article}
\usepackage[utf8]{inputenc}
\usepackage[margin=1in]{geometry}   % 设置页边距
\usepackage{amsmath, amssymb}       % 数学符号与公式
\usepackage{enumitem}               % 自定义列表
\usepackage{unicode-math}
\usepackage[english]{babel}         % 设置整个文档为英文

% 标题与作者信息（可自行修改）
\title{Solution for Chapter 2}
\author{Zimo Guo}
\date{}

\begin{document}

\maketitle

% 使用 \noindent 去掉段落缩进，让题号更清晰
\noindent \textbf{2.1}

\begin{enumerate}[label=(\alph*), left=0pt] % left=0pt 使子题与题号对齐
    \item 
    From standard PAC to 2-oracle PAC: D in standard PAC is any distribution. So for $D_{+}$ and 
    $D_{-}$, P is less than or equal to $\epsilon$. It's PAC-learnable in 2-oracle PAC model 
    obviously.

    \item 
    From 2-oracle PAC to standard PAC: The probability of standard PAC is
    \begin{align*}
        P[y \ne h(x)] &= P[(y = 1 \land h(x) = 0) \lor (y = 0 \land h(x) = 1)] \\
        &\le P[(y = 1 \land h(x) = 0)] + P[(y = 0 \land h(x) = 1)] \\
        &\le 2\epsilon
    \end{align*}

    $\epsilon$ is any positive number, so $P[y \ne h(x)] \le \epsilon$.
\end{enumerate}

\vspace{1em}  % 题目之间的垂直间距

\noindent \textbf{2.4}

This solution doesn't work. For non-concentric cycle, the $h_{s}$ we get by training 
sample $S$ may has a translation from the ideal hypothesis $c_{0}$. If the area of 
$h_{s}$ is equal to $\pi r^2 - \epsilon$, that is, its area is equal to the small
cycle in figure 2.5(b). In the case of figure 2.5(b), $h_{s}$ move to the right top
direction of $c_{0}$. Now three regions both have points locate in $h_{s}$. According
to Gertrude's solution, the error of $h_{s}$ should less than $\epsilon$. But it's 
impossible obviously by the figure 2.5(b).


\vspace{1em}

\noindent \textbf{2.6}

\begin{enumerate}[label=(\alph*), left=0pt] 
    \item 
    Using the same notation in example 2.4, we still set four regions in the rectangle. 
    Now some point will filp, while the sum of the area of the error should still be $\epsilon$.
    Some points in $h_{s}$ will contribute, so every region should have $(\epsilon - \epsilon _{0}) 
    / 4$, where $\epsilon _{0}$ is the error of flipped points. Number of flipped points is less 
    than $m$, so $\epsilon _{0} \le \epsilon \eta$.

    Use the same process in example 2.4, 
    \begin{align*}
        P[R(R_{0}^{'}) > \epsilon] &\le 4 (1 - \frac{\epsilon - \epsilon _{0}}{4}) ^{m} \\
        &= 4 (1 - \frac{\epsilon (1 - \eta)}{4}) ^{m} \\
        &\le 4 \exp \{- \frac{m \epsilon (1 - \eta)}{4} \} \\
        &\le 4 \exp \{- \frac{m \epsilon (1 - \eta ^{'})}{4} \} \\
    \end{align*}

    Let $\delta$ greater than or equal to the right side, get 
    \begin{align*}
        m \ge \frac{4}{\epsilon (1 - \eta ^{'})} \log \frac{4}{\delta}
    \end{align*}
    
\end{enumerate}

\vspace{1em}

\noindent \textbf{2.7}

\begin{enumerate}[label=(\alph*), left=0pt] 
    \item 
    $d(h)$ is the probability that the true label be different from the predict label given
    by $d(h)$. So $d(h ^{*})$ is the probability of filpping, which is $\eta$.

    \item 
    For $R(h)$, the point which has true label 1 will be labelled 0, so $P_{1} = R(h) (1 - \eta)$.
    For area inside $h$, the point which has true label 0 will be labelled 1, so $P_{2} = 
    (1 - R(h)) \eta$. So $d(h) = P_{1} + P_{2} = \eta + (1 - 2 \eta) R(h)$.

    \item 
    \begin{align*}
        d(h) - d(h ^{*}) &= \eta + (1 - 2 \eta) R(h) - \eta \\
        &= (1 - 2 \eta) R(h) \\
        &\ge (1 - 2 \eta) \epsilon \\ 
        &\ge \epsilon (1 - 2 \eta ^{'}) \\
        &= \epsilon ^{'}
    \end{align*}

    \item 
    Let $Z = 1 _{y \ne h ^{*}(x)}$, then $\hat{d}(h^{*}) = (1 / m) \sum_{i=1}^{m} 1 _{y_{i} \ne h
    ^{*} (x_{i})}$. $\hat{d}$ denotes the fraction and $d$ denotes the probability, so $d(h^{*}) = 
    E[\hat{d}(h^{*})]$. Apply Hoeffding's inequality, 

    \begin{align*}
        P[\hat{d}(h^{*}) - d(h^{*}) \ge t] &= P[(1 / m) \sum Z_{i} - E[\hat{d}(h^{*})] \ge t] \\
        &= P[\sum Z_{i} - m E[\hat{d}(h^{*})] \ge mt] \\
        &\le \exp \{-2m t^{2}\} \\
    \end{align*}
    
    Let $t = \frac{\epsilon^{'}}{2}$, 
    \begin{align*}
        P &\le \exp \{-\frac{m \epsilon^{'2}}{2}\} \\
        &\le \frac{\delta}{2} \\
    \end{align*}

    So 
    \begin{align*}
        P[\hat{d}(h^{*}) - d(h^{*}) \ge \frac{\epsilon^{'}}{2}] \le \frac{\delta}{2}
    \end{align*}, that is 

    \begin{align*}
        P[\hat{d}(h^{*}) - d(h^{*}) \le \frac{\epsilon^{'}}{2}] \ge 1 - \frac{\delta}{2} 
    \end{align*}

    \item 
    Similar to (4), we want to prove 
    \begin{align*}
        P[E[\hat{d} (h)] - \frac{1}{m} \sum Z_{i} \ge \frac{\epsilon ^{'}}{2}] \le \frac{\delta}{2}
    \end{align*}

    According to Hoeffding's inequality, 
    \begin{align*}
        left side &= P[\sum Z_{i} - m E[\hat{d} (h)] \le \frac{-m \epsilon ^{'}}{2}] \\
        &\le \exp\{-\frac{m \epsilon^{'2}}{2}\} \\
        &\le \frac{1}{|\mathcal{H}|} \frac{\delta}{2} \\ 
        &\le \frac{\delta}{2}
    \end{align*}

    \item 
    Use the inequality in (3)(4)(5), 

    \begin{align*}
        \hat{d}(h) - \hat{d}(h^{*}) 
        &= [\hat{d}(h) - d(h)] + [d(h) - d(h^{*})] + [d(h^{*}) - \hat{d}(h^{*})] \\ 
        &\ge -\frac{\epsilon^{'}}{2} + \epsilon^{'} - \frac{\epsilon^{'}}{2} \\
        &= 0
    \end{align*}

    So, when $m \ge \frac{2}{\epsilon^{'}} \log \frac{2}{\delta}$, with the probability at least 
    $1 - \frac{\delta}{2}$, $\hat{d}(h) - \hat{d}(h^{*}) \ge 0$ holds.

    Because $\frac{2}{\epsilon^{'}} \log \frac{2}{\delta} \le \frac{2}{\epsilon^{'}} (
    \log |\mathcal{H}| + \log \frac{2}{\delta})$, and $1 - \frac{\delta}{2} > 1 - \delta$, 
    following holds,

    When $m \ge \frac{2}{\epsilon^{'}} (\log |\mathcal{H}| + \log \frac{2}{\delta})$, with the 
    probability at least $1 - \delta$, $\hat{d}(h) - \hat{d}(h^{*}) \ge 0$.
\end{enumerate}

\vspace{1em}

\noindent \textbf{2.12}

According to Hoeffding's inequality, $P[\hat{R_{s}}(h) - R(h) \ge \epsilon] \le \exp \{-2m
\epsilon^{2}\}$. Let $\delta^{'} = p(h)\delta \ge$ the right side of Hoeffding's inequality.
Then we get $\epsilon \ge \sqrt{\frac{\log \frac{1}{p(h)} + \log \frac{1}{\delta}}{2m}}$ with 
probability at least $1 - \delta^{'}$. $1 - \delta \le 1 - \delta^{'}$, so the objective inequality
is true.

\end{document}